\section{Progettare la grafica}

La comunicazione visiva è fondamentale nell'interazione uomo macchina, motivo per cui l'usabilità è fortemente influenzata dall'\bld{interfaccia grafica}. Secondo l'\itc{ISO 9241-12,144}, «la presentazione dell’informazione visiva dovrebbe abilitare l’utente a eseguire i compiti percettivi in modo efficace, efficiente e con soddisfazione».

Durante la progettazione dell'informazione visiva, si devono considerare:

\begin{itemize}
    \item chiarezza
    \item discriminabilità
    \item concisione
    \item consistenza (o coerenza)
    \item scopribilità
    \item leggibilità
    \item comprensibilità
\end{itemize}

\subsection{Leggi della Gestalt}

L'idea portante della Gestalt è che non è corretto dividere l'esperienza umana nelle sue componenti elementari analizzate separatamente, ma bisogna analizzare tutto l'insieme e non solo la somma delle sue parti.

Gli elementi nel campo visivo tendono a organizzarsi in unità secondo la loro forma e posizione relativa:

\begin{itemize}
    \item \bld{Legge della vicinanza}: gli elementi più vicini tra loro tendono a essere raggruppati insieme. Ciò suggerisce di porre uno vicino all'altro elementi grafici che funzionalmente risultano correlati.
    \item \bld{Legge della somiglianza}: gli elementi che più sono simili tra di loro tendono a essere raccolti in unità. Usare forme o colori simili per elementi funzionalmente correlati.
    \item \bld{Legge della chiusura}: la chiusura di una superficie prevale sulla vicinanza degli oggetti. Racchiudere elementi in delle cornici chiuse risulta molto efficace per il loro raggruppamento.
    \item \bld{Legge della continuità di direzione}: le linee che vanno nella stessa direzione tendono a raggrupparsi più facilmente.
    \item \bld{Legge della buona forma}: il campo percettivo segmenta lo spazio in modo tale da far risultare entità equilibrate, armoniche, e coerenti.
    \item \bld{Legge dell'esperienza passata}: gli elementi che danno vita a una figura familiare tendono a formare un'unità.
\end{itemize}
